\chapter{Machine Learning in ROS2 integrieren}

\begin{lernziele}
Nach diesem Kapitel verstehen Sie, wie ein ML-Modell als ROS2-Node in eine Roboterarchitektur eingebunden wird. Sie können Datenfluss, Inferenz, Ergebnisnachrichten, Latenz, Modellqualität und typische Integrationsprobleme beschreiben.
\end{lernziele}

\section{Einleitung}
Machine Learning wird in der Robotik häufig für Wahrnehmung eingesetzt: Objekterkennung, Personenerkennung, Segmentierung oder Klassifikation. Im Beispielprojekt verwenden wir ML zur Personenerkennung.

\section{Rolle von Machine Learning}
Klassische Software arbeitet mit expliziten Regeln. ML-Modelle lernen Muster aus Daten. Für Personenerkennung ist das nützlich, weil Personen in Bildern sehr unterschiedlich aussehen können. Gleichzeitig bringt ML neue Herausforderungen: Trainingsdaten, Fehlklassifikationen, Latenz und schwer erklärbares Verhalten.

\section{Architektur}
\begin{verbatim}
Camera
  -> /camera/image_raw
  -> object_detector_node
  -> /detected_objects
  -> navigation_node oder user_interface_node
\end{verbatim}

Der ML-Node ist Teil der Wahrnehmungsarchitektur. Er sollte klar von Navigation und Sicherheitslogik getrennt werden. Die Navigation bekommt Ergebnisse, aber sie sollte nicht blind jedem Ergebnis vertrauen.

\section{Inference-Node}
Der ML-Node lädt ein trainiertes Modell, verarbeitet eingehende Sensordaten und veröffentlicht Ergebnisse. Dieser Vorgang wird als \emph{Inferenz} bezeichnet: das Anwenden eines bereits trainierten Modells auf neue Daten, im Unterschied zum vorherigen Training des Modells. In Python kann das mit PyTorch, TensorFlow oder ONNX Runtime umgesetzt werden.

\section{Pseudocode}
\begin{lstlisting}[language=Python]
class ObjectDetectorNode(Node):
    def __init__(self):
        super().__init__('object_detector_node')
        self.model = load_model('model.pt')
        self.subscription = self.create_subscription(
            Image,
            '/camera/image_raw',
            self.image_callback,
            10
        )
        self.publisher = self.create_publisher(
            DetectionArray,
            '/detected_objects',
            10
        )

    def image_callback(self, msg):
        image = convert_ros_image_to_tensor(msg)
        detections = self.model(image)
        self.publisher.publish(convert_to_ros_msg(detections))
\end{lstlisting}

Die Dateiendung \texttt{.pt} im Beispiel deutet auf ein PyTorch-Modell hin; bei TensorFlow oder ONNX Runtime als Framework sähen Modell-Dateiformat und Ladefunktion entsprechend anders aus. Auch der Nachrichtentyp \texttt{DetectionArray} ist hier vereinfacht dargestellt: In der Praxis würde man beispielsweise den Standardtyp \texttt{vision\_msgs/Detection2DArray} verwenden oder einen eigenen Nachrichtentyp definieren.

\section{Datenkette}
Die eigentliche Schwierigkeit liegt oft nicht im Modell, sondern in der Datenkette:

\begin{itemize}
    \item Bildformat
    \item Auflösung
    \item Farbraum
    \item Kamerakalibrierung
    \item Latenz
    \item GPU-Nutzung (Grafikprozessor)
    \item Synchronisation
    \item Ergebnisformat
\end{itemize}

Jeder dieser Punkte kann in der Praxis zu Fehlern führen: Ein falscher Farbraum oder eine falsche Auflösung liefert dem Modell unbrauchbare Daten, und eine zu hohe Latenz kann veraltete Erkennungsergebnisse zur Folge haben.

Abbildung~\ref{fig:kap18-datenkette-ml-inferenz} veranschaulicht diese Datenkette vom Kamerabild über die Tensor-Konvertierung und Modell-Inferenz bis zur Veröffentlichung und Weiterverarbeitung der Erkennungsergebnisse.

\begin{figure}[H]
  \centering
  \resizebox{\textwidth}{!}{% =====================================================================
% Abbildung: Datenkette der ML-Inferenz im ObjectDetector
% Passend zu Kapitel 18, Abschnitt "Datenkette"
%
% Einbindung im Buch, z. B. in chapters/18_ml_integration.tex:
%
%   \begin{figure}[H]
%     \centering
%     \resizebox{\textwidth}{!}{% =====================================================================
% Abbildung: Datenkette der ML-Inferenz im ObjectDetector
% Passend zu Kapitel 18, Abschnitt "Datenkette"
%
% Einbindung im Buch, z. B. in chapters/18_ml_integration.tex:
%
%   \begin{figure}[H]
%     \centering
%     \resizebox{\textwidth}{!}{% =====================================================================
% Abbildung: Datenkette der ML-Inferenz im ObjectDetector
% Passend zu Kapitel 18, Abschnitt "Datenkette"
%
% Einbindung im Buch, z. B. in chapters/18_ml_integration.tex:
%
%   \begin{figure}[H]
%     \centering
%     \resizebox{\textwidth}{!}{\input{figures/fig_kap18_datenkette_ml_inferenz}}
%     \caption{Datenkette der ML-Inferenz im \texttt{object\_detector\_node}
%       vom Kamerabild bis zur Veroeffentlichung der Erkennungsergebnisse.}
%     \label{fig:kap18-datenkette-ml-inferenz}
%   \end{figure}
%
% Benoetigte Pakete (im Buch bereits in main.tex geladen): tikz, graphicx
% Benoetigte TikZ-Bibliotheken: positioning, arrows.meta, shadows, calc
% =====================================================================

\input{figures/fig_kap18_datenkette_ml_inferenz.tikz}
}
%     \caption{Datenkette der ML-Inferenz im \texttt{object\_detector\_node}
%       vom Kamerabild bis zur Veroeffentlichung der Erkennungsergebnisse.}
%     \label{fig:kap18-datenkette-ml-inferenz}
%   \end{figure}
%
% Benoetigte Pakete (im Buch bereits in main.tex geladen): tikz, graphicx
% Benoetigte TikZ-Bibliotheken: positioning, arrows.meta, shadows, calc
% =====================================================================

\begin{tikzpicture}[
    scale=0.67, transform shape,
    every node/.style={font=\small},
    node distance=6mm and 6mm,
    ext/.style={
        rectangle, rounded corners=1.5pt, draw=black!50, thick, dashed,
        minimum width=28mm, minimum height=10mm,
        align=center, fill=white, drop shadow
    },
    ros/.style={
        rectangle, rounded corners=1.5pt, draw=orange!60!black, thick,
        minimum width=30mm, minimum height=10mm,
        align=center, fill=orange!15, drop shadow
    },
    code/.style={
        rectangle, rounded corners=1.5pt, draw=orange!60!black, thick,
        minimum width=32mm, minimum height=10mm,
        align=center, fill=orange!8, drop shadow, font=\small\ttfamily
    },
    flow/.style={-{Stealth[length=2mm]}, thick, orange!70!black},
    note/.style={font=\footnotesize, align=center, text width=26mm}
]

% --- Zeile 1 ---
\node[ext] (camera)  at (0, 0)    {Camera\\(Hardware)};
\node[ros] (topic1)  at (4.4, 0)  {Topic\\\texttt{/camera/image\_raw}};
\node[code] (conv1)   at (9.2, 0)  {convert\_ros\_\\image\_to\_tensor()};
\node[ros] (tensor)  at (13.8, 0) {Tensor};

% --- Zeile 2 (darunter, versetzt für Pipeline-Fortsetzung) ---
\node[code] (infer)   at (13.8, -4) {Modell-Inferenz:\\self.model(image)};
\node[ros] (det)     at (9.2, -4)  {Detections\\(roh)};
\node[code] (conv2)   at (4.4, -4)  {convert\_to\_\\ros\_msg()};
\node[ros] (topic2)  at (0, -4)   {Topic\\\texttt{/detected\_objects}};

% --- Zeile 3: Verzweigung ---
\node[ros] (nav) at (-2.6, -7) {navigation\_node};
\node[ros] (ui)  at (2.6, -7)  {user\_interface\_node};

% --- Pfeile Zeile 1 ---
\draw[flow] (camera.east) -- node[note, above, text width=20mm, xshift=-3mm]{Bildformat,\\Auflösung,\\Kalibrierung} (topic1.west);
\draw[flow] (topic1.east) -- (conv1.west);
\draw[flow] (conv1.east)  -- (tensor.west);

% --- Übergang Zeile 1 -> Zeile 2 ---
\draw[flow] (tensor.south) -- node[note, right]{Latenz,\\GPU-Nutzung} (infer.north);

% --- Pfeile Zeile 2 (von rechts nach links) ---
\draw[flow] (infer.west) -- (det.east);
\draw[flow] (det.west)   -- node[note, above]{Synchronisation,\\Ergebnisformat} (conv2.east);
\draw[flow] (conv2.west) -- (topic2.east);

% --- Verzweigung ---
\draw[flow] (topic2.south) -- (nav.north);
\draw[flow] (topic2.south) -- (ui.north);

\end{tikzpicture}

}
%     \caption{Datenkette der ML-Inferenz im \texttt{object\_detector\_node}
%       vom Kamerabild bis zur Veroeffentlichung der Erkennungsergebnisse.}
%     \label{fig:kap18-datenkette-ml-inferenz}
%   \end{figure}
%
% Benoetigte Pakete (im Buch bereits in main.tex geladen): tikz, graphicx
% Benoetigte TikZ-Bibliotheken: positioning, arrows.meta, shadows, calc
% =====================================================================

\begin{tikzpicture}[
    scale=0.67, transform shape,
    every node/.style={font=\small},
    node distance=6mm and 6mm,
    ext/.style={
        rectangle, rounded corners=1.5pt, draw=black!50, thick, dashed,
        minimum width=28mm, minimum height=10mm,
        align=center, fill=white, drop shadow
    },
    ros/.style={
        rectangle, rounded corners=1.5pt, draw=orange!60!black, thick,
        minimum width=30mm, minimum height=10mm,
        align=center, fill=orange!15, drop shadow
    },
    code/.style={
        rectangle, rounded corners=1.5pt, draw=orange!60!black, thick,
        minimum width=32mm, minimum height=10mm,
        align=center, fill=orange!8, drop shadow, font=\small\ttfamily
    },
    flow/.style={-{Stealth[length=2mm]}, thick, orange!70!black},
    note/.style={font=\footnotesize, align=center, text width=26mm}
]

% --- Zeile 1 ---
\node[ext] (camera)  at (0, 0)    {Camera\\(Hardware)};
\node[ros] (topic1)  at (4.4, 0)  {Topic\\\texttt{/camera/image\_raw}};
\node[code] (conv1)   at (9.2, 0)  {convert\_ros\_\\image\_to\_tensor()};
\node[ros] (tensor)  at (13.8, 0) {Tensor};

% --- Zeile 2 (darunter, versetzt für Pipeline-Fortsetzung) ---
\node[code] (infer)   at (13.8, -4) {Modell-Inferenz:\\self.model(image)};
\node[ros] (det)     at (9.2, -4)  {Detections\\(roh)};
\node[code] (conv2)   at (4.4, -4)  {convert\_to\_\\ros\_msg()};
\node[ros] (topic2)  at (0, -4)   {Topic\\\texttt{/detected\_objects}};

% --- Zeile 3: Verzweigung ---
\node[ros] (nav) at (-2.6, -7) {navigation\_node};
\node[ros] (ui)  at (2.6, -7)  {user\_interface\_node};

% --- Pfeile Zeile 1 ---
\draw[flow] (camera.east) -- node[note, above, text width=20mm, xshift=-3mm]{Bildformat,\\Auflösung,\\Kalibrierung} (topic1.west);
\draw[flow] (topic1.east) -- (conv1.west);
\draw[flow] (conv1.east)  -- (tensor.west);

% --- Übergang Zeile 1 -> Zeile 2 ---
\draw[flow] (tensor.south) -- node[note, right]{Latenz,\\GPU-Nutzung} (infer.north);

% --- Pfeile Zeile 2 (von rechts nach links) ---
\draw[flow] (infer.west) -- (det.east);
\draw[flow] (det.west)   -- node[note, above]{Synchronisation,\\Ergebnisformat} (conv2.east);
\draw[flow] (conv2.west) -- (topic2.east);

% --- Verzweigung ---
\draw[flow] (topic2.south) -- (nav.north);
\draw[flow] (topic2.south) -- (ui.north);

\end{tikzpicture}

}
  \caption{Datenkette der ML-Inferenz im \texttt{object\_detector\_node} vom Kamerabild bis zur Veröffentlichung der Erkennungsergebnisse.}
  \label{fig:kap18-datenkette-ml-inferenz}
\end{figure}

\section{Qualität und Sicherheit}
ML-Ergebnisse sind nicht perfekt. Ein Modell kann Personen übersehen oder Gegenstände fälschlich als Personen erkennen. Deshalb sollten Anforderungen und Tests realistisch formuliert werden. Für sicherheitskritische Funktionen sollte ML nicht ohne weitere Sicherheitsmechanismen verwendet werden.

\begin{praxisbox}
Wenn der ML-Node eine Person erkennt, kann die Navigation vorsichtiger fahren. Wenn der ML-Node keine Person erkennt, bedeutet das aber nicht automatisch, dass der Bereich sicher ist. Lidar-basierte Hinderniserkennung bleibt weiterhin wichtig.
\end{praxisbox}

Abbildung~\ref{fig:kap18-ml-sicherheitsarchitektur} stellt diese Kontrastierung zwischen der unsicheren ML-Erkennung und der zuverlässigeren Lidar-basierten Distanzmessung dar.

\begin{figure}[H]
  \centering
  % =====================================================================
% Abbildung: ML-Node in der Sicherheitsarchitektur
% Passend zu Kapitel 18, Abschnitt "Qualität und Sicherheit"
%
% Einbindung im Buch, z. B. in chapters/18_ml_integration.tex:
%
%   \begin{figure}[H]
%     \centering
%     % =====================================================================
% Abbildung: ML-Node in der Sicherheitsarchitektur
% Passend zu Kapitel 18, Abschnitt "Qualität und Sicherheit"
%
% Einbindung im Buch, z. B. in chapters/18_ml_integration.tex:
%
%   \begin{figure}[H]
%     \centering
%     % =====================================================================
% Abbildung: ML-Node in der Sicherheitsarchitektur
% Passend zu Kapitel 18, Abschnitt "Qualität und Sicherheit"
%
% Einbindung im Buch, z. B. in chapters/18_ml_integration.tex:
%
%   \begin{figure}[H]
%     \centering
%     \input{figures/fig_kap18_ml_sicherheitsarchitektur}
%     \caption{Der ML-basierte \texttt{object\_detector\_node} liefert eine
%       unsichere Zusatzinformation, waehrend die Lidar-basierte
%       Hinderniserkennung die zuverlässigere Sicherheitsgrundlage bleibt.}
%     \label{fig:kap18-ml-sicherheitsarchitektur}
%   \end{figure}
%
% Benoetigte Pakete (im Buch bereits in main.tex geladen): tikz
% Benoetigte TikZ-Bibliotheken: positioning, arrows.meta, shadows, calc
% =====================================================================

\input{figures/fig_kap18_ml_sicherheitsarchitektur.tikz}

%     \caption{Der ML-basierte \texttt{object\_detector\_node} liefert eine
%       unsichere Zusatzinformation, waehrend die Lidar-basierte
%       Hinderniserkennung die zuverlässigere Sicherheitsgrundlage bleibt.}
%     \label{fig:kap18-ml-sicherheitsarchitektur}
%   \end{figure}
%
% Benoetigte Pakete (im Buch bereits in main.tex geladen): tikz
% Benoetigte TikZ-Bibliotheken: positioning, arrows.meta, shadows, calc
% =====================================================================

\begin{tikzpicture}[
    every node/.style={font=\small},
    node distance=10mm and 20mm,
    ros/.style={
        rectangle, rounded corners=1.5pt, draw=orange!60!black, thick,
        minimum width=42mm, minimum height=11mm,
        align=center, fill=orange!15, drop shadow
    },
    navnode/.style={ros, minimum width=36mm},
    flowunsafe/.style={-{Stealth[length=2.2mm]}, thick, black!60, dashed},
    flowsafe/.style={-{Stealth[length=2.2mm]}, thick, black!60},
    lbl/.style={font=\tiny, fill=white, inner sep=1pt, align=center}
]

% --- Pfad 1 (oben): ML ---
\node[ros] (mldet) at (0, 3) {\texttt{object\_detector\_node}\\(ML, Kamera)};

% --- Pfad 2 (unten): Lidar ---
\node[ros] (lidar) at (0, -3) {\texttt{lidar\_node} $\rightarrow$\\Hinderniserkennung};

% --- Gemeinsames Ziel ---
\node[navnode] (nav) at (8, 0) {\texttt{navigation\_node}};

\draw[flowunsafe] (mldet.east) -- node[lbl, above, text width=34mm]
    {Person erkannt?\\(unsicher)} (nav.north west);
\draw[flowsafe] (lidar.east) -- node[lbl, below, text width=34mm]
    {physische Distanzmessung\\(zuverlässiger)} (nav.south west);

% --- Hinweistext ---
\node[font=\small, align=center, text width=110mm, below=14mm of nav, xshift=-30mm]
    {ML ergänzt, ersetzt aber nicht die Lidar-basierte Sicherheitsprüfung.};

\end{tikzpicture}


%     \caption{Der ML-basierte \texttt{object\_detector\_node} liefert eine
%       unsichere Zusatzinformation, waehrend die Lidar-basierte
%       Hinderniserkennung die zuverlässigere Sicherheitsgrundlage bleibt.}
%     \label{fig:kap18-ml-sicherheitsarchitektur}
%   \end{figure}
%
% Benoetigte Pakete (im Buch bereits in main.tex geladen): tikz
% Benoetigte TikZ-Bibliotheken: positioning, arrows.meta, shadows, calc
% =====================================================================

\begin{tikzpicture}[
    every node/.style={font=\small},
    node distance=10mm and 20mm,
    ros/.style={
        rectangle, rounded corners=1.5pt, draw=orange!60!black, thick,
        minimum width=42mm, minimum height=11mm,
        align=center, fill=orange!15, drop shadow
    },
    navnode/.style={ros, minimum width=36mm},
    flowunsafe/.style={-{Stealth[length=2.2mm]}, thick, black!60, dashed},
    flowsafe/.style={-{Stealth[length=2.2mm]}, thick, black!60},
    lbl/.style={font=\tiny, fill=white, inner sep=1pt, align=center}
]

% --- Pfad 1 (oben): ML ---
\node[ros] (mldet) at (0, 3) {\texttt{object\_detector\_node}\\(ML, Kamera)};

% --- Pfad 2 (unten): Lidar ---
\node[ros] (lidar) at (0, -3) {\texttt{lidar\_node} $\rightarrow$\\Hinderniserkennung};

% --- Gemeinsames Ziel ---
\node[navnode] (nav) at (8, 0) {\texttt{navigation\_node}};

\draw[flowunsafe] (mldet.east) -- node[lbl, above, text width=34mm]
    {Person erkannt?\\(unsicher)} (nav.north west);
\draw[flowsafe] (lidar.east) -- node[lbl, below, text width=34mm]
    {physische Distanzmessung\\(zuverlässiger)} (nav.south west);

% --- Hinweistext ---
\node[font=\small, align=center, text width=110mm, below=14mm of nav, xshift=-30mm]
    {ML ergänzt, ersetzt aber nicht die Lidar-basierte Sicherheitsprüfung.};

\end{tikzpicture}


  \caption{Der ML-basierte \texttt{object\_detector\_node} liefert eine unsichere Zusatzinformation, während die Lidar-basierte Hinderniserkennung die zuverlässigere Sicherheitsgrundlage bleibt.}
  \label{fig:kap18-ml-sicherheitsarchitektur}
\end{figure}

\section{Optimierung}
Für reale Roboter können ONNX Runtime, TensorRT oder OpenVINO sinnvoll sein – plattform- beziehungsweise hardwareoptimierte Laufzeitumgebungen, die trainierte Modelle schneller und ressourcenschonender ausführen. Ein großes PyTorch-Modell direkt auf einem kleinen Rechner kann zu langsam sein. Optimierung betrifft nicht nur Geschwindigkeit, sondern auch Speicherverbrauch, Temperatur und Energie.

\section{Modellierung in SysML}
Im SysML-Modell sollte sichtbar sein:

\begin{itemize}
    \item welche Anforderung der ML-Node unterstützt
    \item welche Daten er benötigt
    \item welches Ergebnis er liefert
    \item welche Latenz akzeptabel ist
    \item welcher Test die Erkennung prüft
    \item welche Fehlerfälle berücksichtigt werden
\end{itemize}

\section{Teststrategie für ML-Funktionen}
ML-Funktionen müssen anders getestet werden als klassische Regeln. Ein einzelner Test mit einem Bild reicht nicht aus. Sinnvoll ist ein kleiner Testdatensatz mit unterschiedlichen Lichtverhältnissen, Perspektiven und Situationen. Für den Einstieg kann dieser Datensatz klein sein, aber er sollte bewusst zusammengestellt werden.

Mögliche Testfragen:

\begin{itemize}
    \item Erkennt das Modell Personen bei normaler Beleuchtung?
    \item Wie verhält es sich bei Gegenlicht?
    \item Gibt es viele falsche Erkennungen?
    \item Wie lange dauert die Inferenz pro Bild?
    \item Was passiert, wenn keine Kamera-Nachricht ankommt?
\end{itemize}

Im SysML-Modell kann ein Testfall \texttt{TC-004 Personenerkennung} auf mehrere konkrete Testdatensätze verweisen. Damit bleibt sichtbar, dass die Anforderung nicht nur durch Code, sondern durch Daten und Bewertungskriterien geprüft wird.

\section{Zusammenfassung}
ML kann die Wahrnehmung des Roboters verbessern, bringt aber neue Integrations- und Qualitätsfragen. Ein ML-Node sollte sauber in die ROS2-Architektur eingebunden, im SysML-Modell nachvollziehbar beschrieben und durch Tests überprüft werden.

\section{Kontrollfragen}
\begin{enumerate}
    \item Wofür wird ML im Roboterprojekt eingesetzt?
    \item Welche Daten benötigt der ObjectDetector?
    \item Welches Topic könnte Erkennungsergebnisse veröffentlichen?
    \item Warum ist Latenz wichtig?
    \item Warum sollte ML nicht blind als Sicherheitsfunktion betrachtet werden?
    \item Welche Optimierungswege gibt es für ML-Inferenz?
\end{enumerate}

\section{Übungen}
\begin{uebungbox}
\textbf{Übung 1: Modellbezug}\\
Beschreiben Sie im SysML-Modell, welche Anforderung durch den ML-Node erfüllt wird und welches Topic die Ergebnisse veröffentlicht.
\end{uebungbox}

\begin{uebungbox}
\textbf{Übung 2: Fehlerfälle}\\
Nennen Sie drei mögliche Fehlerfälle bei der Personenerkennung und beschreiben Sie eine sinnvolle Systemreaktion.
\end{uebungbox}

\section{Literaturhinweise}
Für ROS2-Integration ist die ROS2-Dokumentation \cite{ros2docs} wichtig. Für ML-Modelle sollten zusätzlich Quellen zu PyTorch, YOLO, ONNX und Embedded-AI-Optimierung genutzt werden.
